% =============================================================================
% Ready-to-paste LaTeX: emotion regression (Experiment 1) + cross-dataset emotion
% transfer to Blockbuster, with the PCA-8 bottleneck control. Uses bib keys present in
% library.bib (kang2025there, kim2010music, ma2021computational, choi2017transfer,
% eerola2011comparison). Numbers from experiments/emotion/exp_emotion_regression.py,
% exp_emotion_improve.py and exp_blockbuster_emotion.py (RandomForest regressor).
%
% PROTOCOL: the Eerola column of tab:emotion_transfer comes from the corrected run
% (results/cv_corrected.json: 10x5 repeated GroupKFold, nested-CV-tuned C, macro-F1 on
% pooled out-of-fold predictions, training emotions predicted out-of-fold) so that it
% agrees with statistical_power.tex. Do NOT reintroduce the per-fold figures
% (0.394 / 0.397 / 0.343 / 0.350 / 0.166, from results/statistical_power.json), which are
% biased low (progress log section 21), nor the older single-five-fold figures
% (0.406 / 0.388 / 0.350 / 0.328), which came from one partition at the sklearn-default C.
% The 5-genre random-guess floor (0.309) is the mean genre prevalence, which is what a
% base-rate random guess scores in expectation; where a random guess was measured, the
% formula agrees to within 0.005 (6-genre Eerola 0.251 vs 0.250; Blockbuster 0.297 vs
% 0.292/0.295). The BLOCKBUSTER column is likewise now the repeated-CV,
% nested-CV-tuned run with per-cue emotion bridging (results/blockbuster_deep.json,
% progress log section 17); do not reintroduce the older single-five-fold figures there
% either (0.565 / 0.559 / 0.582).
% =============================================================================

\subsection{Predicting Emotion from Audio Embeddings}
\label{subsec:emotion_regression}

The first pipeline stage predicts the eight emotion ratings from an audio representation.
Two choices were examined: the regression model and the feature set. For the model, ridge
regression, support vector regression with an RBF kernel, and a random forest were compared
on the VGGish features; the choice proved consequential, with mean $R^2$ over the eight
emotions rising from $0.37$ (ridge) to $0.54$ (SVR) and $0.56$ (random forest). The
non-linear random forest was therefore adopted. Under this regressor the three feature sets
-- the Audio Spectrogram Transformer embedding (AST, 768-d), the VGGish embedding (128-d) and
a librosa hand-crafted MIR set (103-d) -- were compared by GroupKFold cross-validation
grouped by film (Table~\ref{tab:emotion_regression}).

\begin{table}[t]
  \centering
  \caption{Emotion regression: mean $R^2$ over the eight emotions (random forest, GroupKFold
    by film, $n=360$ clips).}
  \label{tab:emotion_regression}
  \begin{tabular}{lr}
    \hline
    Feature set & mean $R^2$ \\
    \hline
    AST (768) & 0.560 \\
    VGGish (128) & 0.558 \\
    MIR / librosa (103) & 0.490 \\
    \hline
  \end{tabular}
\end{table}

With a non-linear regressor the AST and VGGish embeddings are essentially tied, and the
hand-crafted MIR features are only slightly behind. A mean $R^2$ of $0.56$ is within the
range reported for music emotion recognition on comparable data \citep{kang2025there}. As is
standard in that literature, the arousal-related dimensions (energy, tension, valence) are
recovered most reliably, whereas the discrete positive emotions -- happiness and sadness --
remain the hardest to predict, reflecting the long-standing ``valence problem''
\citep{kim2010music,kang2025there}. VGGish is retained for the transfer experiment below not
because it is the strongest emotion predictor but because it is the only representation that
can be reproduced identically for both datasets.

\subsection{Cross-Dataset Emotion Transfer and a Bottleneck Control}
\label{subsec:emotion_transfer}

The Blockbuster dataset of \citet{ma2021computational} provides genre labels and
pre-extracted VGGish features but no audio and no emotion annotations, so it cannot by itself
support an emotion pipeline. Because VGGish is a fixed pretrained model, extracting it for the
Eerola clips yields a feature space identical to the one shipped with Blockbuster. A
VGGish$\to$emotion random forest was trained on the entire Eerola corpus and used to predict
emotions for both datasets, from which genre was then classified. To test whether the
\emph{emotion} bottleneck carries information beyond that of an arbitrary low-dimensional
projection, an eight-component PCA of the VGGish features was included as a control.

\begin{table}[t]
  \centering
  \caption{Genre classification through the predicted-emotion bottleneck versus the raw
    embedding and an eight-dimensional PCA control (logistic regression, Macro-F1). The
    Both columns are $10\times5$ repeated cross-validation with nested-CV-tuned
    regularisation. \textbf{The two columns are not comparable to one another} --- they
    differ in label space (5 vs 6 genres), in unit (a 10\,s clip vs a whole film) and
    in chance level (a base-rate random guess scores $0.309$ on the left and $0.295$ on
    the right). Compare \emph{within} a column only; the like-for-like cross-corpus
    comparison is Table~\ref{tab:zero_shot}.}
  \label{tab:emotion_transfer}
  \begin{tabular}{lrr}
    \hline
    Approach & Eerola (5 genres) & Blockbuster (6 genres) \\
    \hline
    predicted emotion $\rightarrow$ genre & \textbf{0.417} & \textbf{0.616} \\
    ground-truth emotion $\rightarrow$ genre & 0.428 & --- \\
    PCA-8 $\rightarrow$ genre (control) & 0.372 & 0.587 \\
    raw VGGish $\rightarrow$ genre & 0.377 & 0.593 \\
    full hand-crafted MIR (140) $\rightarrow$ genre & --- & 0.585 \\
    most-frequent baseline & 0.168 & 0.035 \\
    base-rate random guess & 0.309 & 0.295 \\
    \hline
  \end{tabular}
\end{table}

Table~\ref{tab:emotion_transfer} yields three results. First, on the primary Eerola corpus
the predicted-emotion bottleneck ($0.417$) exceeds both the PCA-8 control ($0.372$) and the
raw embedding ($0.377$): the eight emotions retain more genre-relevant signal than an
arbitrary eight-dimensional compression, so the representation is not merely an interpretable
relabelling of a generic bottleneck. The control itself is no better than the embedding it
compresses ($-0.005$, $p=0.77$), so compression alone does not explain the gain. The margin
over the control is consistent, holding in all ten repeats, and significant under the film
bootstrap ($+0.045$, $p=0.029$) but borderline under the stricter Nadeau--Bengio test
($p=0.122$, Section~\ref{subsec:statistical_power}), and it is reported as such. On the easier,
film-level Blockbuster task the emotion bottleneck leads the control nominally but not
significantly ($0.616$ vs $0.587$, $p=0.369$), so it is reported there as competitive
rather than superior. Notably, the full 140-feature hand-crafted MIR set ($0.585$) is
statistically indistinguishable from the VGGish embedding ($0.593$, $p=0.787$) on that
corpus, so no claim of a learned-embedding advantage is made. Second, the predicted emotions are as genre-informative as the human
ratings ($0.417$ against a ground-truth ceiling of $0.428$, a difference that is not
significant, $p=0.33$), so the regressor recovers the genre-relevant emotional content
without requiring manual annotation. Nor does adding the raw embedding back help: on the
eight-genre task, a naive fusion of the emotion features with the AST embedding scores
$0.260$ against $0.299$ for the emotion features alone. Third,
the transferred emotions are face-valid: on Blockbuster, the films of a genre deviate from
the mean predicted emotion exactly as film-scoring convention predicts -- horror toward fear
($+0.99$), romance toward tenderness ($+1.05$), comedy toward happiness, drama toward
tenderness and action toward anger -- even though the regressor was trained on a different
corpus and never saw Blockbuster. This is a form of the transfer behaviour expected of
pretrained audio embeddings \citep{choi2017transfer}.

Two caveats temper these results. Even over fifty folds, the margins by which the emotion
bottleneck exceeds the control and the direct embeddings on Eerola (of the order of
$0.04$--$0.05$ macro-$F_1$) are significant under the film bootstrap but only borderline
under the Nadeau--Bengio test, whose attainable $p$-value on this corpus lies above $0.05$
for these comparisons; settling them requires more films rather than more computation. And, because Blockbuster contains no human
emotion ratings, the predicted emotions there can be validated only indirectly, through their
downstream genre-predictive value and face validity, not against a ground truth
\citep{ma2021computational,eerola2011comparison}.
